\documentclass[11pt]{article}
% ===========================================================================
%  Shared preamble for the AP Calculus AB exam-preparation worksheets.
%  Input right after \documentclass. Files beginning with "_" are not compiled
%  on their own. Navy and gold, to match the course's lesson pages.
% ===========================================================================
\usepackage[letterpaper,top=2.3cm,bottom=2.2cm,left=1.7cm,right=1.7cm,headheight=28pt,footskip=30pt]{geometry}
\usepackage{amsmath,amssymb}
\usepackage[T1]{fontenc}
\usepackage[utf8]{inputenc}
\usepackage{lmodern}
\usepackage{xcolor}
\usepackage[most]{tcolorbox}
\usepackage{enumitem}
\usepackage{fancyhdr}
\usepackage{array}
\usepackage{tikz}
\usepackage{pgfplots}
\pgfplotsset{compat=1.18}

\definecolor{navy}{HTML}{161E54}
\definecolor{ink}{HTML}{27357E}
\definecolor{gold}{HTML}{B98A00}
\definecolor{goldbg}{HTML}{FFFDF3}
\definecolor{inkbg}{HTML}{F6F7FD}
\definecolor{rule}{HTML}{C9CFF0}
\definecolor{okgreen}{HTML}{1B7A44}
\definecolor{okbg}{HTML}{F1F8F3}

\pagestyle{fancy}
\fancyhf{}
\fancyhead[L]{\footnotesize Name: \underline{\hspace{4cm}}}
\fancyhead[C]{\textbf{\color{navy}Integration Academy}\\[-2pt]{\scriptsize Calculus AB \textperiodcentered{} AP\textsuperscript{\textregistered} Exam Preparation}}
\fancyfoot[L]{\tiny\color{black!55}\parbox[t]{0.78\textwidth}{AP\textsuperscript{\textregistered} is a registered trademark of the College Board, which was not involved in the production of, and does not endorse, this material. \textcopyright{} Integration Academy.}}
\fancyfoot[R]{\scriptsize\thepage}
\renewcommand{\headrulewidth}{1.2pt}
\renewcommand{\headrule}{\hbox to\headwidth{\color{gold}\leaders\hrule height \headrulewidth\hfill}}
\renewcommand{\footrulewidth}{0pt}
\setlength{\parindent}{0pt}
\setlength{\parskip}{4pt}

% Title block: \wstitle{code}{title}{unit line}
\newcommand{\wstitle}[3]{\begin{center}{\small\bfseries\color{gold}WORKSHEET #1}\\[3pt]{\LARGE\bfseries\color{navy}#2}\\[3pt]{\small\color{black!60}#3}\end{center}\vspace{2pt}}
% Section heading
\newcommand{\sech}[1]{\par\vspace{10pt}{\large\bfseries\color{navy}#1}\par\nobreak\vspace{2pt}{\color{rule}\hrule height0.8pt}\vspace{7pt}}
% A level tag printed beside a group of questions
\newcommand{\level}[1]{\par\vspace{6pt}{\footnotesize\bfseries\color{gold}\MakeUppercase{#1}}\par\nobreak\vspace{2pt}}

\newtcolorbox{ideabox}{colframe=ink,colback=inkbg,coltitle=white,colbacktitle=ink,fonttitle=\bfseries\small,title={Key ideas},arc=2pt,boxrule=0.8pt,left=8pt,right=8pt}
\newtcolorbox{example}[1]{colframe=rule,colback=white,coltitle=navy,colbacktitle=inkbg,fonttitle=\bfseries\small,title={#1},arc=2pt,boxrule=0.8pt,borderline west={3pt}{0pt}{ink},left=10pt,right=8pt}
\newtcolorbox{frbox}[1]{colframe=gold,colback=goldbg,coltitle=white,colbacktitle=gold,fonttitle=\bfseries\small,title={#1},arc=2pt,boxrule=0.8pt,breakable,left=8pt,right=8pt}
\newtcolorbox{solbox}{colframe=okgreen,colback=okbg,coltitle=white,colbacktitle=okgreen,fonttitle=\bfseries\small,title={Worked solutions},arc=2pt,boxrule=0.8pt,breakable,left=8pt,right=8pt}
\newtcolorbox{scorebox}[1]{colframe=navy,colback=white,coltitle=white,colbacktitle=navy,fonttitle=\bfseries\small,title={#1},arc=2pt,boxrule=0.8pt,breakable,left=8pt,right=8pt}

\newcommand{\soln}{\par\smallskip\textbf{\color{ink}Solution.}\ }
\newcommand{\result}[1]{\par\smallskip\textbf{#1}}
% blank working space of a given height
\newcommand{\work}[1]{\par\vspace{#1}}
% a point value in the margin of a scoring guide
\newcommand{\pts}[1]{\hfill{\footnotesize\bfseries\color{gold}[#1]}}

% Coordinate grid for graph questions: \begin{apgraph}{xmin}{xmax}{ymin}{ymax} ... \end{apgraph}
\newenvironment{apgraph}[4]{\begin{center}\begin{tikzpicture}\begin{axis}[
  axis lines=middle, xmin=#1, xmax=#2, ymin=#3, ymax=#4,
  grid=both, grid style={black!12}, xtick distance=1, ytick distance=1,
  ticklabel style={font=\scriptsize}, xlabel={$x$}, ylabel={$y$},
  width=10cm, height=7.2cm, samples=80, clip=false,
  every axis plot/.append style={line width=1.1pt, color=ink}]}{\end{axis}\end{tikzpicture}\end{center}}
\newcommand{\closeddot}[2]{\addplot[only marks, mark=*, mark size=2.2pt, color=ink] coordinates {(#1,#2)};}
\newcommand{\opendot}[2]{\addplot[only marks, mark=*, mark size=2.2pt, color=ink, mark options={fill=white}] coordinates {(#1,#2)};}

% Multiple-choice options: four across, or stacked when the options are long.
\newcommand{\opts}[4]{\par\vspace{3pt}\begin{tabular}{@{}p{0.235\linewidth}p{0.235\linewidth}p{0.235\linewidth}p{0.235\linewidth}@{}}(A)\ #1 & (B)\ #2 & (C)\ #3 & (D)\ #4\end{tabular}\par\vspace{9pt}}
\newcommand{\optsv}[4]{\par\vspace{3pt}\begin{tabular}{@{}l@{\ \ }p{0.9\linewidth}@{}}(A) & #1\\(B) & #2\\(C) & #3\\(D) & #4\end{tabular}\par\vspace{9pt}}

\fancyhead[R]{\footnotesize\color{black!60}Worksheet 5.5}
\begin{document}

\wstitle{5.5}{Absolute Extrema on a Closed Interval}{Unit 5: Analysing Functions}

\begin{ideabox}
\begin{itemize}[leftmargin=*,itemsep=2pt,topsep=2pt]
  \item \textbf{Candidates test.} On a closed interval, the absolute extrema occur at critical points or at endpoints.
  \item Evaluate $f$ at every critical point in the interval and at both endpoints. The largest value is the absolute maximum; the smallest is the absolute minimum.
  \item The list of values \emph{is} the justification. Show it.
  \item On an open interval with a single critical point, a sign change of $f'$ over the whole interval gives an absolute extremum.
\end{itemize}
\end{ideabox}

\sech{Worked examples}
\begin{example}{Example 1 \textperiodcentered{} The candidates test}
Find the absolute extrema of $f(x)=x^3-3x^2+1$ on $[-1,4]$.\soln $f'(x)=3x(x-2)$. $f(-1)=-3$, $f(0)=1$, $f(2)=-3$, $f(4)=17$.\result{Maximum 17 (at $x=4$); minimum $-3$ (at $x=-1$ and $x=2$).}
\end{example}

\begin{example}{Example 2 \textperiodcentered{} A trigonometric function}
Find the absolute extrema of $f(x)=\sin x+\dfrac x2$ on $[0,2\pi]$.\soln $f'(x)=\cos x+\dfrac12=0$ at $\dfrac{2\pi}{3}$ and $\dfrac{4\pi}{3}$. $f(0)=0$, $f\!\left(\dfrac{2\pi}{3}\right)\approx1.913$, $f\!\left(\dfrac{4\pi}{3}\right)\approx1.228$, $f(2\pi)=\pi$.\result{Maximum $\pi$ (at $x=2\pi$); minimum 0 (at $x=0$).}
\end{example}

\begin{example}{Example 3 \textperiodcentered{} An open interval}
Find the absolute minimum of $f(x)=x^2+\dfrac{54}{x}$ for $x>0$.\soln $f'(x)=2x-\dfrac{54}{x^2}$ is negative for $0<x<3$ and positive for $x>3$.\result{The absolute minimum is $f(3)=27$.}
\end{example}

\newpage
\sech{Practice}

\level{Core}
\begin{enumerate}[leftmargin=*,itemsep=2pt]
  \item Find the absolute extrema of $f(x)=x^2-2x$ on $[0,3]$.\work{2.3cm}
  \item Find the absolute extrema of $f(x)=4-x^2$ on $[-1,2]$.\work{2.3cm}
  \item Find the absolute extrema of $f(x)=x^3$ on $[-2,1]$.\work{2.3cm}
  \item Why must the endpoints be included in the candidates test?\work{2.3cm}
\end{enumerate}

\level{Exam level}
\begin{enumerate}[leftmargin=*,itemsep=2pt,start=5]
  \item Find the absolute extrema of $f(x)=2x^3-9x^2+12x$ on $[0,3]$.\work{2.8cm}
  \item Find the absolute extrema of $f(x)=xe^{-x/2}$ on $[0,6]$.\work{2.8cm}
  \item Find the absolute extrema of $f(x)=x^2-8\ln x$ on $[1,e]$.\work{2.8cm}
  \item Find the absolute extrema of $f(x)=|x-2|$ on $[0,5]$.\work{2.8cm}
\end{enumerate}

\begin{frbox}{Question 9 \textperiodcentered{} Free-response practice}
Let $f(x)=3x^4-4x^3-12x^2+2$.
\begin{enumerate}[label=(\alph*),leftmargin=*,itemsep=2pt]
  \item Find the critical points of $f$.\work{2.6cm}
  \item Find the absolute maximum value of $f$ on $[-2,3]$. Justify your answer.\work{2.6cm}
  \item Find the absolute minimum value of $f$ on $[-2,3]$.\work{2.6cm}
\end{enumerate}
\end{frbox}

\newpage
\begin{solbox}
\begin{enumerate}[leftmargin=*,itemsep=4pt]
  \item $f(0)=0$, $f(1)=-1$, $f(3)=3$: maximum 3, minimum $-1$.
  \item $f(-1)=3$, $f(0)=4$, $f(2)=0$: maximum 4, minimum 0.
  \item $f(-2)=-8$, $f(0)=0$, $f(1)=1$: maximum 1, minimum $-8$.
  \item An extremum on a closed interval can occur at an endpoint, where the derivative need not be zero.
  \item $f'(x)=6(x-1)(x-2)$. $f(0)=0$, $f(1)=5$, $f(2)=4$, $f(3)=9$: maximum 9, minimum 0.
  \item $f'(x)=e^{-x/2}\left(1-\dfrac x2\right)=0$ at $x=2$. $f(0)=0$, $f(2)=\dfrac2e$, $f(6)=\dfrac{6}{e^3}$: maximum $\dfrac2e$, minimum 0.
  \item $f'(x)=2x-\dfrac8x=0$ at $x=2$. $f(1)=1$, $f(2)=4-8\ln2\approx-1.545$, $f(e)=e^2-8\approx-0.611$: maximum 1, minimum $4-8\ln2$.
  \item The only critical point is $x=2$, where $f'$ does not exist. $f(0)=2$, $f(2)=0$, $f(5)=3$: maximum 3, minimum 0.
  \item \begin{enumerate}[label=(\alph*),leftmargin=*,itemsep=3pt]
    \item $f'(x)=12x^3-12x^2-24x=12x(x-2)(x+1)$: $x=-1$, $x=0$ and $x=2$.
    \item $f(-2)=34$, $f(-1)=-3$, $f(0)=2$, $f(2)=-30$, $f(3)=29$. The absolute maximum is 34, at the endpoint $x=-2$.
    \item From the same list, the absolute minimum is $-30$, at $x=2$.
  \end{enumerate}
\end{enumerate}
\end{solbox}

\end{document}
